\chapter{勾股定理}\label{ch:g8:pythagoras}

数学中最著名的定理，刻画了直角三角形的三条边
之间的关系。借助它，许多无法直接用尺子量出的长度，
例如对角线、距离和高，都能算出来。这也是我们
第一次接触定理与其
逆定理之间的区别。

\section{定理}

\begin{theorem}[毕达哥拉斯]\label{thm:g8:pythagoras:direct}
若三角形 $ABC$ 的直角顶点为 $A$，则
\[
BC^2 = AB^2 + AC^2 :
\]
斜边（直角所对的边）的平方
等于另外两条边的平方和。
\end{theorem}

\begin{proof}[证明思路]
四个全等的三角形拼放在边长为
$AB + AC$ 的大正方形内，留出一个以斜边为边的倾斜正方形，
比较面积即可得出等式。详细计算
见\cref{exo:g8:pythagoras:11}，该定理也可由
高中分册中学习的数量积推出。
\end{proof}

\begin{omfigure}
\begin{tikzpicture}[scale=0.58]
  \coordinate (A) at (0,0);
  \coordinate (B) at (4,0);
  \coordinate (C) at (0,3);
  \draw[thick, fill=omDef!12] (A) -- (B) -- (C) -- cycle;
  \draw[thin] (0.35,0) -- (0.35,0.35) -- (0,0.35);
  \node[below left, font=\small] at (A) {$A$};
  \node[below right, font=\small] at (B) {$B$};
  \node[above left, font=\small] at (C) {$C$};
  % squares on the sides
  \draw[omDef, thick, fill=omDef!8] (0,0) rectangle (4,-4);
  \node[omDef, font=\small] at (2,-2) {$AB^2 = 16$};
  \draw[omMeth, thick, fill=omMeth!8] (0,0) rectangle (-3,3);
  \node[omMeth, font=\small] at (-1.5,1.5) {$AC^2 = 9$};
  \draw[omThm, thick, fill=omThm!8] (4,0) -- (7,4) -- (3,7) -- (0,3)
    -- cycle;
  \node[omThm, font=\small] at (3.5,3.5) {$BC^2 = 25$};
\end{tikzpicture}

{\small 定理的面积表述：以斜边为边的正方形
（$25$）的面积恰好等于另外两个正方形面积之和
（$16 + 9$）。这里三边分别为 $3$、$4$、$5$。}
\end{omfigure}

\begin{method}[求边长]\label{met:g8:pythagoras:compute}
在直角三角形中，已知两条边时：
\begin{enumerate}
  \item 指出斜边（直角所对的边 --- 总是
        最长的边）；
  \item 写出\cref{thm:g8:pythagoras:direct}中的等式，
        并代入已知的值；
  \item 求出未知的平方：斜边未知时，\emph{相加}
        两个平方；直角边未知时，\emph{相减}；
  \item 开平方（可精确计算，也可用计算器），
        并检验斜边是否最长。
\end{enumerate}
\end{method}

\begin{example}[求斜边]\label{ex:g8:pythagoras:hypotenuse}
直角三角形的两条直角边分别为 $6$ 和 $8$，斜边为 $c$：
\[
c^2 = 6^2 + 8^2 = 36 + 64 = 100,
\qquad
c = \sqrt{100} = 10 .
\]
（数 $\sqrt{100}$，即 $100$ 的\emph{平方根}，是
平方为 $100$ 的正数；平方根在
\cref{ch:g9:sqrt}中专门讨论。）
\end{example}

\begin{example}[求直角边]\label{ex:g8:pythagoras:leg}
一架长 $5$ 米的梯子斜靠在墙上，梯脚距墙 $1.4$ 米。
设梯顶的高度为 $h$：梯子即为斜边，于是
\[
h^2 = 5^2 - 1.4^2 = 25 - 1.96 = 23.04,
\qquad
h = \sqrt{23.04} = 4.8 \text{ 米}.
\]
这里应当相减，而不是相加：所求的是\emph{直角边}。
\end{example}

\section{逆定理}

\begin{theorem}[毕达哥拉斯定理的逆定理]\label{thm:g8:pythagoras:converse}
若三角形 $ABC$ 的三边满足 $BC^2 = AB^2 + AC^2$，则
该三角形的直角顶点为 $A$。
\end{theorem}
\admitted

\begin{method}[判断直角三角形]\label{met:g8:pythagoras:test}
已知三角形的三条边：
\begin{enumerate}
  \item 分别计算最长边的平方，以及另外两条边的平方
        之和；
  \item 若两个结果相等，则三角形为直角三角形
        （逆定理），直角顶点位于最长边所对的顶点。
  \item 若两者不相等，则三角形不是直角三角形 --- 因为
        若为直角三角形，毕达哥拉斯定理就要求二者相等。
\end{enumerate}
\end{method}

\begin{example}\label{ex:g8:pythagoras:test}
三边为 $5$、$12$、$13$：最长边的平方 $13^2 = 169$；另外两边
的平方和 $5^2 + 12^2 = 25 + 144 = 169$。两者相等，故为直角三角形
（直角位于边 $13$ 所对的顶点）。

三边为 $6$、$7$、$9$：$9^2 = 81$，而 $6^2 + 7^2 = 85 \neq 81$，
故不是直角三角形（它看上去“差不多”是，
但应由数字而非眼睛来判定）。
\end{example}

\begin{remark}[定理与逆定理]
定理与其逆定理的内容并不相同：定理
\emph{利用}直角来计算长度；逆定理则
利用长度来\emph{证明}直角。工匠们很早便已使用
逆定理，已有数千年之久：一根在 $3$、$4$、$5$ 个单位处打结的绳子，
拉紧后便得到一个标准的直角。
\end{remark}

\begin{example}[正方形的对角线]\label{ex:g8:pythagoras:diagonal}
边长为 $1$ 的正方形，其对角线是
直角边为 $1$、$1$ 的直角三角形的斜边：
\[
d^2 = 1^2 + 1^2 = 2 ,
\]
故 $d = \sqrt2 \approx 1.414$，这是一个
不是分数的数，高中分册将证明这一点。若正方形边长为 $c$，则其对角线为
$c\sqrt2$。
\end{example}

\section{练习}

\begin{exercise}[$\star$]\label{exo:g8:pythagoras:1}
在下列每个直角三角形中，求斜边长：
直角边分别为 $9$ 和 $12$；\quad 直角边为 $5$ 和 $12$；\quad 直角边为 $8$ 和 $15$。
\end{exercise}

\begin{exercise}[$\star$]\label{exo:g8:pythagoras:2}
一个直角三角形的斜边为 $25$，一条直角边为 $7$，求
另一条直角边。
\end{exercise}

\begin{exercise}[$\star$]\label{exo:g8:pythagoras:3}
求两边长分别为 $12$ cm 和 $9$ cm 的长方形的对角线；求边长为 $5$ cm 的
正方形的对角线（先给出含平方根的精确值，
再四舍五入到毫米）。
\end{exercise}

\begin{exercise}[$\star$]\label{exo:g8:pythagoras:4}
哪些三角形是直角三角形？请利用
\cref{met:g8:pythagoras:test}加以说明：
\[
(20,\ 21,\ 29); \qquad
(7,\ 8,\ 11); \qquad
(1.5,\ 2,\ 2.5).
\]
\end{exercise}

\begin{exercise}[$\star$]\label{exo:g8:pythagoras:5}
一扇门由一块斜撑木板加固。门宽 $3$ m、
高 $1.6$ m。木板长多少？
\end{exercise}

\begin{exercise}[$\star$]\label{exo:g8:pythagoras:6}
一个等腰三角形的两腰长均为 $10$ cm，底边长 $12$ cm。
其高将它分成两个直角三角形。求该高，
再求三角形的面积。
\end{exercise}

\begin{exercise}[$\star\star$]\label{exo:g8:pythagoras:7}
一台电视屏幕是长宽比为 $16{:}9$ 的长方形，其尺寸为 $88.5$ cm 乘
$49.8$ cm。电视屏幕按对角线长（以英寸计）出售，
（$1$ 英寸 $= 2.54$ cm）。求对角线的长（以 cm 为单位），再
求以英寸计的长度（四舍五入到最接近的整数英寸）。
\end{exercise}

\begin{exercise}[$\star\star$]\label{exo:g8:pythagoras:8}
一艘船先向东航行 $24$ km，再向北航行 $10$ km。问它离出发点
的直线距离是多少？请先画出示意图。
\end{exercise}

\begin{exercise}[$\star\star$]\label{exo:g8:pythagoras:9}
一架 $2.5$ m 长的梯子须能够到离地面 $2.4$ m 高的窗台。
梯脚应距墙多远？所得结果是否
符合安全建议（梯脚距墙约为梯长的
四分之一）？
\end{exercise}

\begin{exercise}[$\star\star$]\label{exo:g8:pythagoras:10}
在坐标网格中描出点 $A(1, 2)$ 与 $B(5, 5)$，并画出一个
直角边平行于坐标轴的直角三角形，其斜边为
$[AB]$。求 $AB$。（对任意两点而言，这一作图
将引出解析几何中的距离公式，
其内容见高中分册。）
\end{exercise}

\begin{exercise}[$\star\star\star$]\label{exo:g8:pythagoras:11}
四个直角边为 $a$、$b$、斜边为 $c$ 的相同直角三角形，
分别放在边长为 $a + b$ 的正方形的四个角上，中间留出一个
边长为 $c$ 的倾斜正方形。
\begin{enumerate}
  \item 用两种方式表示大正方形的面积：一是直接计算，
        二是写成（四个三角形）$+$（倾斜正方形）。
  \item 展开 $(a + b)^2$（见\cref{thm:g8:equations:expand}），
        并由此推出 $c^2 = a^2 + b^2$：你便证明了
        毕达哥拉斯定理。
\end{enumerate}
\end{exercise}

\section{问题：勾股数}

\begin{problem}[{周末问题 --- 整数直角三角形，
以及恒等式 $(m^2 - n^2)^2 + (2mn)^2 = (m^2 + n^2)^2$，
它能生成全部此类三角形}]
\label{pb:g8:pythagoras:1}
称\emph{勾股数}\index{勾股数}为一组整数
$(a, b, c)$，满足
\[
a^2 + b^2 = c^2 ,
\]
如工匠们所用的 $(3, 4, 5)$。古巴比伦的一块泥板记录了
更大的一些勾股数，其中包括 $(4601,\ 4800,\ 6649)$，
比毕达哥拉斯早一千多年。本题寻找勾股数，
然后构造一台能够产生它们的“机器”：一个代数恒等式，
正是\cref{ch:g8:equations}中的内容，被用于
服务于几何学。

\medskip
\noindent\textbf{第一部分 --- 寻找勾股数。}
\begin{enumerate}
  \item 验证 $(5, 12, 13)$、$(8, 15, 17)$ 和
        $(20, 21, 29)$ 都是勾股数。
        由\cref{thm:g8:pythagoras:converse}，
        边长为这些数的三角形具有何种性质？
  \item 说明为何将一根由十二段等长线段组成、拉紧的绳子
        摆成边长分别为 $3$、$4$、$5$ 段的三角形，
        便能为工匠提供标准的直角。
  \item 证明：若 $(a, b, c)$ 是勾股数，
        则 $(ka, kb, kc)$ 对任意整数 $k \geq 1$ 也是勾股数。
        由此从 $(3, 4, 5)$ 得出三组新的勾股数。
  \item 所谓 $(3, 4, 5)$ 的倍数，即三元组
        $(3k, 4k, 5k)$。证明 $(5, 12, 13)$ 并
        \emph{不是}其中之一。
  \item 说明为何第 3 题与第 4 题合起来可以证明
        对 $(3, 4, 5)$ 作缩放永远不可能得到所有勾股数，
        必须借助另一台“机器”。
\end{enumerate}

\medskip
\noindent\textbf{第二部分 --- 机器。}
取两个整数 $m > n \geq 1$，构造三个数
\[
a = m^2 - n^2,
\qquad
b = 2mn,
\qquad
c = m^2 + n^2 .
\]
\begin{enumerate}[resume]
  \item 利用
        \cref{pb:g8:equations:1}中的乘法公式（记 $A = m^2$、$B = n^2$），
        展开 $a^2$ 与 $b^2$，并证明
        \[
        \left(m^2 - n^2\right)^2 + (2mn)^2
        = \left(m^2 + n^2\right)^2 :
        \]
        这台“机器”总是输出一组勾股数。
  \item 将“机器”作用于 $(m, n) = (2, 1)$、$(3, 1)$、$(3, 2)$、
        $(4, 1)$ 与 $(4, 3)$，并将五组三元组
        列成一张表。
  \item 表中有一组三元组是另一组较小三元组的缩放。
        是哪一组？它由哪一组缩放得到？
  \item 求 $(m, n)$，使“机器”输出第 1 题中的
        三元组 $(20, 21, 29)$（两直角边次序不论）。
  \item 求 $m$ 与 $n$，使 $m^2 + n^2 = 100$，
        并由此得出斜边为 $100$ 的一组勾股数。
\end{enumerate}

\medskip
\noindent\textbf{第三部分 --- 勾股数族与奇偶性总结。}
\begin{enumerate}[resume]
  \item 取 $m = n + 1$，令“机器”运行。证明：奇数直角边为
        $2n + 1$，偶数直角边为 $2n^2 + 2n$，斜边为
        $2n^2 + 2n + 1$ --- 因此斜边与偶数直角边相差
        恰好为 $1$。分别写出 $n = 1, 2, 3, 4$ 时的三元组。
  \item 由此推出：\emph{每一个}奇数 $3, 5, 7, 9, \dots$
        都是某个整数直角三角形的一条边，并给出
        一组包含 $11$ 的三元组。
  \item 偶数同样可行：令 $n = 1$ 运行“机器”，并
        得到一组包含 $14$ 的三元组。然后，
        对第 7 题表中的某一组作缩放，
        再次得到同一组三元组。
  \item 证明：偶数的平方是偶数，奇数的平方是奇数。
        （将该数写成 $2k$ 或 $2k + 1$，
        并利用一个恒等式。）
  \item 作为奇偶性总结：不存在三个数都是奇数的勾股数 ---
        即在每个整数直角三角形中，
        至少有一条边为偶数。
\end{enumerate}
\end{problem}
